\section{Extremal positive transport and its aspect}
\label{sec:classical-extremality}

Fix $0<\gamma<1$, and put
\[
 w(b,a)=(1+\gamma)b+(1-\gamma)a
       =a+b-\gamma(a-b),\qquad 0\le b\le a.
\]
All entropies in this section are divided by $\log(1/\eps)$.
The particle $T$ is a random variable carrying two bounded vectors
$V(T),A(T)\in\RR^2$; its position at time $c$ is $X(c)=A+cV$.
The bounded real variable $c$ may have any joint law with $T$.
The observations used below are finite.  For a finite label $Y$ and
arbitrary $T$, $\Ent(Y\mid T)$ denotes the average entropy of the
conditional probabilities $\EE[\ind_{\{Y=y\}}\mid\sigma(T)]$, divided
by $\log(1/\eps)$, and $\Mut(T;Y)=\Ent(Y)-\Ent(Y\mid T)$.  Thus no entropy
of $T$ is used.  A fixed enlargement of the bounded domains is allowed.
There is no Planck parameter in this section.

\subsection{Apertures and the universal growth rate}

Use nested grids for time and write $C_t$ for the cell of $c$ of length
comparable to $\eps^t$, together with its chosen representative.  At $t=0$
one cell contains the time domain and its representative is zero.  The
finitely many boundary choices in this convention have entropy
$O(1/\log(1/\eps))$.  For an unoriented normal $R$, let $Q_R$ be an
orthogonal frame whose first vector is tangent and whose second vector is
normal.  Orientations of aspect $e$ are recorded in cells of angular size
comparable to $\eps^e$.  Nested angular grids, or a finite collection of
nested slope grids, may be used.  Choose a representative normal for each
recorded angular cell; both rectangular grids use its frame.

\begin{definition}[Standard aperture observation]
\label{def:classical-aperture}
At duration $t$ and horizon $L$, a standard observation $G$ records
\[
 (t,b,a,C_t,R_{a-b},[V]_{b,a;R},
                   [X(C_t)]_{b+t,a+t;R}),
 \qquad 0\le b\le a,\quad a+t\le L.
\]
Here the bracket denotes the rectangular grid cell, in the recorded
frame, with side lengths $\eps^b,\eps^a$, or respectively
$\eps^{b+t},\eps^{a+t}$.  Its aperture weight is $w(G)=w(b,a)$.
The shape and orientation may be chosen by a stochastic kernel; extra
finite labels are permitted.  Equivalently, a label may locate the data
in bounded enlargements of the indicated rectangles.  Adding the actual
grid cells then costs $O(1/\log(1/\eps))$ information.
\end{definition}

Depths can be rounded on a mesh $\xi=\xi_\eps\downarrow0$ with
$\log(1/\xi)=o(\log(1/\eps))$.  The number of shape choices in a bounded
triangle is then subpower.  In this terminology a factor is subpower if
its normalized logarithm tends to zero.  Dyadic spatial widths give a
finite menu of apertures: choose orientation spacing comparable to the
minor-to-major axis ratio and translates spaced by the corresponding
widths.  Every rectangle is contained in a fixed enlargement of one in
this menu.  Indeed an angular error at most a constant times the axis
ratio moves a tangent uncertainty into at most a constant multiple of
the normal uncertainty.  This verifies the assertion for both columns.
Discarding extra labels never increases particle information.  Define the
root supremum by
\[
 \tau_L(T)=\sup_{D\text{ at duration }0}
       \{\EE w(D)-\Mut(T;D)\}.
\]
The supremum is taken for the current particle law and the current
precision, before any limit.  For a fixed finite menu it is a supremum
over all stochastic assignments to compatible caps.  Rounding depths down
on the mesh and adjoining the actual target cells costs
$O(\xi)+O(1/\log(1/\eps))$ in root score, uniformly on fixed domain,
horizon, and enlargement bounds.  Thus these menus recover the unrestricted
root supremum in the limit.

The score balances three quantities. The aperture weight rewards finer
localization, particle information charges what the observation reveals
about $T$, and conditional time entropy measures the time information
left after the particle is known. The root supremum provides the
zero-duration benchmark for this balance.

Write $S_t=\Ent(C_t\mid T)$ and
\begin{equation}
 \mathcal F_t(G)=\EE w(G)-\Mut(T;G)+2S_t,
 \qquad
 \betac=\sup\limsup_{\eps\to0}
       \frac{\mathcal F_t(G)-\tau_L(T)}{t}.
 \label{eq:classical-score}
\end{equation}
The supremum ranges over all fixed $0<t\le L<\infty$ and all sequences
of the systems just defined.  Domain bounds are fixed within each
sequence.  Fixed affine normalizations put them in one bounded domain
without changing a limit.  The label $G$ determines, up to bounded
ambiguity, a root aperture of the same shape: recover $A=X(C_t)-C_tV$
and coarsen its position rectangle.  Consequently
$\mathcal F_t(G)-\tau_L(T)\le2S_t+o(1)\le2t+o(1)$.
For a single particle and independent uniform time, the choices
$b=a=L-t$ give $\tau_L=2L+o(1)$ and $\mathcal F_t=2L+o(1)$.
Thus
\begin{equation}
 0\le\betac\le2.
 \label{eq:classical-provisional}
\end{equation}

\begin{theorem}[Positive-transport growth]
\label{thm:classical-growth}
For the aperture systems of Definition~\ref{def:classical-aperture},
the rate in \eqref{eq:classical-score} satisfies $\betac\le\gamma$.
\end{theorem}

We prove this bound by examining a sequence whose growth approaches
$\betac$. Localization and the least-aspect construction produce a
saturated path. Its tangent reduction has two alternatives: an isotropic
system, excluded in Section~\ref{sec:isotropic}, or a path whose aspect
grows in proportion to time. Sections~\ref{sec:characteristic-plane}
and~\ref{sec:rate-transfer} constrain the local information rates of
the latter path; Section~\ref{sec:classical-closure} combines those
constraints with saturation to finish the proof. Until that point we
use only \eqref{eq:classical-provisional} and the definition of $\betac$.

The rate in \eqref{eq:classical-score} is unchanged if the underlying
probability spaces are required to be finite.  Fix the complete finite
root menu and project $G$ to its geometric coordinates, which can only
decrease particle information.  For every possible output/time pair $(g,k)$,
use the two compatibility tests for $V$ and $A+\bar c_kV$, where
$\bar c_k$ is the recorded time representative.  Partition particles
by all these tests and by membership in every root cap.  The finite
signature $\kappa(T)$ depends on $T$ alone, not on the sampled time.
Choose an actual $(V,A)$ in each positive signature, with finite
particle index $T^{\mathrm f}$.  Give $(T^{\mathrm f},G,C_t)$ the old
law of $(\kappa(T),G,C_t)$, and realize each time label by a point in
its actual cell.  The resulting finite law is compatible and has the
same expected weight, with
\[
 \Mut(T^{\mathrm f};G)\le\Mut(T;G),\qquad
 \Ent(C_t\mid T^{\mathrm f})\ge\Ent(C_t\mid T).
\]
For roots, averaging a compatible kernel within signatures preserves
its output marginal and weight and decreases information; conversely,
every signature kernel lifts through $\kappa(T)$ with equal information.
The frozen-menu root suprema therefore agree exactly.  The uniform menu
error above is $o(1)$, so every general-law score excess is at most that
of a finite law plus $o(1)$.  We take these finite laws before forming
the conditional systems and independent copies below.

We make explicit the uniformity implicit in that definition.  For any
fixed domain bound $B$, $K<\infty$, and $t_{\min}>0$, there is a function
$\omega_\eps(B,K,t_{\min})\to0$ such that
\begin{equation}
 \mathcal F_t(G)\le\tau_L(T)+\betac t+
 \omega_\eps(B,K,t_{\min}),\qquad t_{\min}\le t\le L\le K.
 \label{eq:classical-uniform}
\end{equation}
Here the menu, its depth mesh, and its fixed enlargement convention are
chosen as deterministic functions of $\eps$, common to all systems in
the supremum.  It also applies to every conditional system with these
bounds.  To see
this, a failure selects systems and parameters with a fixed positive
excess and a convergent subsequence of $(t,L)$.  Changing depths by $\nu$
costs at most $C_{K,\gamma}\nu+o(1)$: cover the larger rectangle by the
finer cells, including its changed frame in the count.  Coarsening time
by $\nu$ loses at most $\nu+o(1)$ of conditional time entropy.  Transport
to the coarser time uses an uncertainty $\eps^{t-\nu}$ times the measured
velocity rectangle, which fits the earlier position rectangle.  Finally,
coarsening both depths of a root test by at most $\nu$ loses at most
$2\nu+o(1)$ in weight; this proves the same continuity for the root
supremum when its horizon changes by $\nu$.  Fixed parameter slack
therefore converts the selected systems to a sequence contradicting
\eqref{eq:classical-score}.  The same argument at $t=0$ uses the root
supremum directly, with no division by time.

Whenever we use conditional arrays below, we first fix the finite list
of observations, the depth mesh, all trimming thresholds, and the
positive time increments.  Their exceptional probabilities and
normalized errors are then bounded by one deterministic function
$q_\eps\to0$ on every retained cell.  Cell coordinate bounds and
covering multiplicities are common constants.  Thus every selection
of one retained system per precision remains an admissible sequence
in \eqref{eq:classical-score}.  The conditional-array selection principle at the end of
Section~\ref{sec:entropy} applies to precisely this array; it does
not require a uniform rate over unspecified subpower envelopes.

\subsection{Splitting and cap bounds}

Call an observation saturated if, along the sequence under discussion,
$\mathcal F_t(G)-\tau_L(T)\to\betac t$.  An extremizer below always means
such a sequence, not a maximizing measure at finite precision.

\begin{lemma}[Splitting saturated laws]
\label{lem:classical-splitting}
Let $Z$ be a finite tag.  If $\Delta=\betac t+\tau_L(T)-\mathcal F_t(G)$,
and $\Delta_z$ is the corresponding deficit for the law conditional on
$Z=z$, then
\[
 \EE\Delta_Z\le\Delta+3\Ent(Z).
\]
In particular subpower splittings preserve saturation on typical parts,
simultaneously for any fixed finite list of saturated observations.
\end{lemma}
\begin{proof}
Combining conditional root tests and including $Z$ in their label gives
$\EE\tau_L(T\mid Z)\le\tau_L(T)+\Mut(T;Z)$.
The exact chain rules give
\[
 \EE\mathcal F_t(G\mid Z)-\mathcal F_t(G)
 =-\Mut(T;Z\mid G)+\Mut(T;Z)-2\Mut(C_t;Z\mid T).
\]
Subtract and use that each of the last two nonnegative costs is at most
$\Ent(Z)$.  By \eqref{eq:classical-uniform}, each conditional deficit is
at least $-\omega_\eps$, uniformly over the cells.  Its positive part
therefore tends to zero in mean when the original deficit and tag
entropy do.  Markov's inequality proves the last assertion.  No lower
bound on an individual cell probability is used for this selection.
\end{proof}

\begin{lemma}[Greedy cap regularization]
\label{lem:classical-cap-regularization}
After a subpower splitting and removal of probability $o(1)$, a particle
law can be restricted to classes with a bounded number $h$ such that
\begin{equation}
 \Pr_T(\operatorname{cap}(D))
       \le\eps^{w(D)-h-o(1)},\qquad \tau_L(T)=h+o(1).
 \label{eq:classical-cap-bounds}
\end{equation}
The cap bound is uniform over the root menu and its fixed enlargements.
For a further law $Q\le\eps^{-\eta}\Pr_T$ on such a class,
$\tau_L(Q)\le h+\eta+o(1)$.  Hence a subpower restriction or tilt
preserving the displayed score profiles preserves saturation.
\end{lemma}
\begin{proof}
In the finite root menu, repeatedly select a cap maximizing
$w(D)+\logeps{\mu_{\rm rem}(\operatorname{cap}(D))}$ and remove its
remaining contents.  Selected contents $E_i$ are disjoint and cover the
law.  A selected cap cannot be selected again with positive mass, so the
procedure is finite.  The successive maxima $h_i$ decrease.  Split the
indices into intervals $[h,h+\xi]$ and keep each entire $E_i$ in its bin.
At the first removal in this bin every cap has remaining mass at most
$\eps^{w(D)-h-\xi}$.  The absolute mass of its intersection with the
whole bin has the same upper bound.

Let the bin have probability $p$.  Under its normalized law every cap
thus has mass at most $p^{-1}\eps^{w(D)-h-\xi}$.  The deterministic
assigned-cap test has score
\[
 \sum_{i\text{ in bin}}\frac{\mu(E_i)}p
       \left(w(D_i)+\logeps{\frac{\mu(E_i)}p}\right)
 =\EE[h_i\mid\text{bin}]+\logeps{1/p}.
\]
Conversely, for any compatible stochastic test $D$, conditional relative
entropy on its cap gives
\[
 \Mut(T;D)\ge
   \EE\logeps{\frac1{\Pr_T(\operatorname{cap}(D))}}.
\]
The two estimates put the conditional root supremum between
$h+\logeps{1/p}$ and $h+\xi+\logeps{1/p}$, up to the menu ambiguity.
The same inequality proves the assertion about $Q$.

For completeness, very negative bins can be discarded: a fixed number
of zero-depth caps cover the domain, so a remaining maximum below
$-\lambda$ forces remaining probability at most $C\eps^\lambda$.
The other maxima are at most $2L+o(1)$.  Choose the mesh slowly, and
discard bins of mass less than $\eps^{\eta_\eps}$, where
$\eta_\eps\downarrow0$ slowly enough that the number of bins times
$\eps^{\eta_\eps}$ tends to zero.  Then $\logeps{1/p}=o(1)$ on every
retained bin.  This proves \eqref{eq:classical-cap-bounds}.  Apply
Lemma~\ref{lem:classical-splitting} when the input is saturated.  Before
any later subpower tilt, impose the uniform probability bins of
Lemma~\ref{lem:homogenization}; its conditional relative-entropy
estimate preserves the stipulated profiles.  Their score tends to
$\betac t+h$, while the cap bound gives $\tau_L(Q)\le h+o(1)$ and
\eqref{eq:classical-uniform} gives the reverse inequality.
\end{proof}

\subsection{Anisotropic localization}

\begin{lemma}[Localization and composition]
\label{lem:classical-localization}
Let $G$ have fixed shape $(b,a)$ at $t$, put $e=a-b$, and fix
$0\le r<t$, $0\le d\le e$.  Let $J$ be the aperture at $r$ of shape
$(b,b+d)$ whose frame is the depth-$d$ prefix of the frame of $G$.
Set $\zeta=\Ent(R_d\mid T,C_r)$.  Conditional on $J$, there are bounded
line systems with final shape $(0,e-d)$, duration $t-r$, and horizon
$L'=e-d+t-r$.  With $\tau'_J$ their root suprema,
\begin{align}
 \mathcal F_t(G)
 &\le\betac(t-r)+w(b,b+d)-\Mut(T;J)\notag\\
 &\qquad{}+2S_r+\EE\tau'_J+2\zeta+o(1),
 \label{eq:classical-local-first}\\
 w(b,b+d)-\Mut(T;J)+2S_r+\EE\tau'_J
 &\le\betac r+\tau_L(T)+o(1).
 \label{eq:classical-local-second}
\end{align}
Here parameters are fixed before $\eps\to0$.  Uniform versions have
errors $C\xi+C\log(C_0)/\log(1/\eps)+\omega_\eps$ for depth-rounding
mesh $\xi$, fixed enlargement bound $C_0$, and the compact parameter
bounds in \eqref{eq:classical-uniform}.  If $\zeta=o(1)$ and $G$ is
saturated, typical conditional systems are saturated and near-maximizing
compositions give a saturated predecessor at $r>0$.
\end{lemma}
\begin{proof}
Write $v_J,x_J$ for the velocity and position representatives of $J$,
$c_J$ for its time representative, and
$S=Q_{R_d}\operatorname{diag}(\eps^b,\eps^{b+d})$.  The coordinates
\[
 v'=S^{-1}(V-v_J),\quad
 a'=\eps^{-r}S^{-1}(A+c_JV-x_J),\quad
 c'=\eps^{-r}(c-c_J)
\]
are bounded on this conditional system, and its trajectory is
$a'+c'v'$.  More explicitly,
\[
 a'+c'v'=\eps^{-r}S^{-1}
       \bigl(A+cV-x_J-(c-c_J)v_J\bigr),
\]
so the moving anchor being subtracted is known from $J$.  Both columns
use the same spatial map; the position column
has the additional factor $\eps^{-r}$.  If $\theta$ is the angle from
the coarse frame to the final frame, then
$|\sin\theta|\le C\eps^d$.  In these units the final velocity
uncertainty matrix, with its harmless scalar factor removed, is
\[
 N=\operatorname{diag}(1,\eps^{-d})
       Q_\theta\operatorname{diag}(1,\eps^e).
\]
Its largest singular value lies between positive fixed constants:
the first column has length at least one and all entries are bounded.
Its determinant has absolute value $\eps^{e-d}$.  Its other singular
value is therefore comparable to $\eps^{e-d}$.  The final position
matrix is $\eps^{t-r}N$.  This proves the asserted final aperture.

The same calculation proves the needed determination of $J$ by $G$.
Rotate the velocity uncertainty to the coarse frame, coarsen its normal
width from $a$ to $b+d$, and transport the position uncertainty from
$C_t$ to $C_r$.  Since $|C_t-C_r|\le C\eps^r$, the transported velocity
uncertainty fits the position widths $\eps^{b+r},\eps^{b+d+r}$.
Thus only a bounded number of cells is possible.  This argument compares
uncertainty rectangles, and imposes no bound on the displacement of
their centers.  In particular
\[
 \Ent(J\mid G)=o(1),\qquad
 \Ent(J\mid T,C_r)\le\zeta+o(1).
\]
After adding the actual transformed cells at bounded cost, the final
label $G'$ satisfies
$\Mut(T;J)+\Mut(T;G'\mid J)\le\Mut(T;G)+o(1)$.
Also, since $J$ includes $C_r$,
\[
 0\le S_t-S_r-\Ent(C_t\mid T,J)
      =\Mut(C_t;J\mid T,C_r)\le\zeta+o(1).
\]
The local final weight is exactly
$w(0,e-d)=w(b,a)-w(b,b+d)$.  Apply
\eqref{eq:classical-uniform} in each conditional system and average;
these identities give \eqref{eq:classical-local-first}.

For composition, take a local root test of shape $(B,A)$,
$0\le B\le A\le L'$, and frame $Q_\phi$.  Its global uncertainty is
\[
 M=S Q_\phi\operatorname{diag}(\eps^B,\eps^A).
\]
Let $p\le q$ be its two singular-value depths.  The determinant identity,
the upper bound for the largest singular value, the lower bound for the
smallest one, and submultiplicativity of the condition number give
\begin{gather*}
 p+q=2b+d+A+B+o(1),\qquad
 q-p\le d+A-B+o(1),\\
 p\ge b+B-o(1),\qquad
 q\le b+d+A+o(1)\le a+t-r+o(1).
\end{gather*}
Consequently
\[
 w(p,q)\ge w(b,b+d)+w(B,A)-o(1),\qquad q+r\le L+o(1).
\]
The position column has the same singular directions and the extra
factor $\eps^r$, so this is an admissible global aperture at $r$.
Rounding its depths down removes the displayed horizon slack at a
vanishing cost.  Include $J$ in the composed label and add its actual
cells at bounded ambiguity.  Its information is at most
$\Mut(T;J)+\Mut(T;D'\mid J)+o(1)$.
Choose conditional root tests within any prescribed $\eta>0$ of
$\tau'_J$, apply the global bound at $r$, and then let $\eta\downarrow0$.
At $r=0$ use the defining root supremum instead.  This proves
\eqref{eq:classical-local-second}.

Finally, if both endpoints of these inequalities agree in the limit,
their nonnegative limiting deficits have zero mean.  The uniform lower
bounds on all conditional deficits permit the same Markov selection as
in Lemma~\ref{lem:classical-splitting}.  The composition of the chosen
conditional root tests gives the predecessor assertion.  The finite
errors quoted in the statement come respectively from rounding,
bounded covering multiplicities, and \eqref{eq:classical-uniform}.
\end{proof}



The case $r=d=0$ removes the common depth $b$.  Replacing the local
horizon by its smallest permitted value $e+t$ cannot lower a score
minus its root supremum; the universal upper bound then forces
saturation.  Replacing $\eps$ by $\eps^t$ makes the duration one and
divides all depths, weights, and normalized entropies by $t$.

\subsection{The least aspect and saturated predecessors}

We first isolate the rectangle comparison used twice below.  Two
equal-shape apertures $(b,b+e)$ with angular separation of depth
$v\in[0,e]$ have a common coarsening of shape $(b,b+v)$ and an
intersection aperture of shape $(b+e-v,b+e)$.  The common label is
determined by either original label, once an angle-depth bin of width
$\xi$ is supplied, at information cost $O(\xi)+o(1)$.  For the
intersection, one normal constraint has width $\eps^{b+e}$; solving the
other normal constraint along its tangent gives width
$C\eps^{b+e}/|\sin\theta|$.  The safe endpoint of the angle bin gives
the asserted tangent depth with an $O(\xi)$ loss.  Coarsening to normal
depth $b+v$ proves the common-rectangle assertion.  These arguments
apply separately to both columns, with the common time factor when
the duration is positive.  The weight gain is
\begin{equation}
 w(b+e-v,b+e)+w(b,b+v)-2w(b,b+e)=2\gamma(e-v).
 \label{eq:classical-rectangle-gain}
\end{equation}

Figure~\ref{fig:aperture-geometry} shows the two rectangles responsible
for this weight gain.
\begin{figure}[H]
\centering
\begin{tikzpicture}[x=0.91cm,y=0.91cm,font=\small,
  aperture/.style={line width=0.7pt},
  comparison/.style={draw=figuregray,dashed,line width=0.6pt},
  dimension/.style={<->,>=Stealth,line width=0.45pt}]
  \begin{scope}[shift={(3.05,0)}]
    \node[font=\small\bfseries] at (0,2.10) {Common coarsening};
    \fill[figuregray!7] (-2.8,-1.35) rectangle (2.8,1.35);
    \draw[comparison] (-2.8,-1.35) rectangle (2.8,1.35);
    \draw[aperture] (-2.7,-0.18) rectangle (2.7,0.18);
    \begin{scope}[rotate=25]
      \draw[aperture,figureblue] (-2.7,-0.18) rectangle (2.7,0.18);
    \end{scope}
    \node[anchor=south west] at (2.12,0.19) {$D_1$};
    \node[anchor=south east,text=figureblue] at (2.35,1.22) {$D_2$};
    \draw[line width=0.45pt] (0.88,0) arc[start angle=0,end angle=25,radius=0.88];
    \node[anchor=west,font=\footnotesize,fill=white,inner sep=1pt]
      at (0.82,0.49) {$\theta\asymp\eps^v$};
    \draw[dimension] (-2.8,-1.65) -- (2.8,-1.65)
      node[midway,below=2pt] {$\eps^b$};
    \draw[dimension] (-3.07,-1.35) -- (-3.07,1.35)
      node[midway,above=2pt,sloped] {$\eps^{b+v}$};
    \node[font=\footnotesize] at (0,-2.65) {depths $(b,b+v)$};
  \end{scope}
  \begin{scope}[shift={(11.0,0)}]
    \node[font=\small\bfseries] at (0,2.10) {Intersection};
    \draw[aperture] (-2.7,-0.18) rectangle (2.7,0.18);
    \begin{scope}[rotate=25]
      \draw[aperture,figureblue] (-2.7,-0.18) rectangle (2.7,0.18);
    \end{scope}
    \fill[figureblue!22]
      (-0.81193,-0.18) -- (0.03991,-0.18) --
      (0.81193,0.18) -- (-0.03991,0.18) -- cycle;
    \draw[comparison] (-0.81193,-0.18) rectangle (0.81193,0.18);
    \node[anchor=south west] at (2.12,0.19) {$D_1$};
    \node[anchor=south east,text=figureblue] at (2.35,1.22) {$D_2$};
    \draw[dimension] (-0.81193,-0.69) -- (0.81193,-0.69)
      node[midway,below=2pt] {$\eps^{a-v}$};
    \draw[dimension] (1.13,-0.18) -- (1.13,0.18);
    \node[anchor=west,fill=white,inner sep=1pt] at (1.26,0) {$\eps^a$};
    \node[font=\footnotesize,align=center] at (0,-1.65)
      {each original aperture:\quad $\eps^b\times\eps^a$};
    \node[font=\footnotesize] at (0,-2.65) {depths $(a-v,a)$};
  \end{scope}
\end{tikzpicture}
\caption{Aperture geometry for two common-shape observations.
The centered schematic has $0<\eps<1$, $b\le a$ and
$0\le v\le a-b$, with angle $\theta\asymp\eps^v$.
Here $a=b+e$ in the notation of \eqref{eq:classical-rectangle-gain};
the two comparison rectangles give its weight gain.
The dashed common rectangle contains both apertures up to fixed factors;
the shaded intersection is contained in a rectangle with depths
$(a-v,a)$.  All indicated side lengths are orders of magnitude, and the
picture displays an interior value of $v$.  The same calculation applies
to the position column at duration $r$ after multiplying both side
lengths by $\eps^r$.}
\label{fig:aperture-geometry}
\end{figure}


\begin{proposition}[Least-aspect limiting extremizer]
\label{prop:least-aspect}
There exists a saturated sequence at duration one with $b=0$, endpoint
aspect $e=m$, and horizon $1+m$, where $m\ge0$ is the least aspect of
any such saturated sequence.  One may moreover require $Y_j$ conditional on $X_j$ to be homogeneous
for a prescribed countable list of finite pairs $(X_j,Y_j)$, provided
that for each fixed $j$ there is $0\le A_j<\infty$, independent of $\eps$,
such that $|\mathcal Y_j|\le\eps^{-A_j}$ for all sufficiently small $\eps$.
No cardinality bound is imposed on $X_j$, which may include the exact
particle $T$.
\end{proposition}
\begin{proof}
Choose systems whose normalized rates tend to $\betac$, split by narrow
shape bins, remove $b$ by localization, and normalize time.  Each
individual inner sequence now has a fixed finite aspect $e$; its
rate defect can be made arbitrarily small.  We need a bound for $e$
independent of which sequence was selected.

Sample two final labels independently given $T$, and let $D_1,D_2$ be
their root coarsenings of shape $(0,e)$.  If the rate defect is $\eta$,
then, since $S_1\le1+o(1)$ and $\betac\ge0$,
\[
 \Mut(T;D_j)\le(1-\gamma)e-\tau_L(T)+2+\eta+o(1).
\]
Tag their separation depth $v$ on a fixed mesh and use the common and
intersection labels above.  Conditional independence gives
\begin{align*}
 2e-(1+\gamma)\EE v-\tau_L(T)
 &\le\Mut(T;D_1,D_2)+O(\xi)+o(1)\\
 &\le2\big((1-\gamma)e-\tau_L(T)+2+\eta\big)
       -\big((1-\gamma)\EE v-\tau_L(T)\big)
       +O(\xi)+o(1).
\end{align*}
Indeed the information identity for the pair subtracts
$\Mut(D_1;D_2)$, which is at least the common-label entropy minus its
two determination errors.  That entropy is at least its particle
information, and the root test bounds the latter below by its weight
minus $\tau_L$.  The separation tag has entropy $o(1)$ at each fixed
mesh.  Rearrangement yields
\[
 2\gamma\EE(e-v)\le4+2\eta+O(\xi)+o(1).
\]
Thus a fixed $K=K_\gamma<\infty$ gives a positive, uniformly bounded
below probability of agreement within a bounded number of cells at
depth $(e-K)_+$.  For each $T$, maximize the conditional probability
of a coarse orientation cell.  The conditional collision probability
of two independent draws is at most a fixed constant times this
maximum.  This deterministic predictor therefore succeeds, to bounded
cell ambiguity, on an event of probability at least $c_\gamma>0$.
On that event the coarse orientation has conditional entropy $o(1)$
given $T$.  Binary splitting preserves a rate defect tending to zero
on this event.  Localization with $r=0$, $d=(e-K)_+$ produces a
conditional system with residual aspect at most $K$ and rate defect
tending to zero.  Uniform conditional deficit bounds justify selecting
one such system even when the conditioning cells are numerous.

It remains to take the least aspect.  The preceding construction,
followed by a diagonal, produces at least one saturated sequence with
bounded aspect and reduced horizon.  Let $m$ be the infimum of its
possible limiting aspects.  Choose such sequences with aspects
$e_j\to m$.  For the $j$th sequence choose its precision small enough
that its score defect, finite-menu errors, and the errors for the first
$j$ requested profiles are less than $1/j$.  Coarsen an endpoint of
aspect $e_j$ to aspect $m$ when $e_j\ge m$; if needed allow a depth error
$|e_j-m|$ before this coarsening.  The score loses at most
$C_\gamma|e_j-m|+o(1)$, whereas reducing the horizon to $1+m$ reduces
the subtracted root supremum.  The universal bound forces the resulting
diagonal to saturate.  For the homogeneity assertion, apply Lemma~\ref{lem:homogenization} to
finite initial lists of these pairs under the stated bounds on the child
alphabets, and choose their meshes and errors before the corresponding
precision.  Lemma~\ref{lem:profile-diagonal} then retains the countable
list.  In the applications below the children are finite tuples of grid
observations at fixed depths in bounded coordinates; the exact particle
and other unrestricted conditioning data occur only among the parents.
\end{proof}

\begin{proposition}[Saturated path and coherence]
\label{prop:saturated-path}
The sequence in Proposition~\ref{prop:least-aspect} can carry, on one
joint law of $(T,c)$, saturated observations $G_u$ at every positive
dyadic time $u\le1$.  Their shapes satisfy
\[
 b(1)=0,\quad a(1)=m,\quad
 b(r)\ge b(t),\quad a(r)+r\le a(t)+t\quad(r<t),
 \qquad e(u):=a(u)-b(u)\ge mu.
\]
They can be constructed so that
$\Ent(R_{e(u)}\mid T,C_u)=o(1)$.  There is a constant $K_\gamma$,
independent of $r,t$, for which the orientations at $r<t$ agree with
probability tending to one at every positive depth
\begin{equation}
 d<\min(e(r),e(t))-K_\gamma
       \bigl(|b(r)-b(t)|+|a(r)-a(t)|+t-r\bigr).
 \label{eq:classical-coherence}
\end{equation}
All statements refer to the inner $\eps$ limit with the times fixed.
\end{proposition}
\begin{proof}
On a finite dyadic mesh, work backwards from $1$ using
Lemma~\ref{lem:classical-localization} with $d=0$.  A composed predecessor
has minimum depth at least the later $b$ and maximum depth at most
the later $a$ plus the time difference.  Split its variable shape into
narrow bins and use Lemma~\ref{lem:classical-splitting} to preserve all
equalities already imposed.  Drop redundant labels.  Before making
this split, sample each incoming predecessor from its kernel given
$(T,C_r)$, independently of the previously retained finer time and
later labels.  This operation preserves its law with $(T,C_r)$ and
preserves the entire later joint law.  In particular the underlying
time variable is never resampled.

Subsequent finite binning and trimming introduce at most their normalized
entropy or relative-entropy cost into the averaged conditional mutual
information, by Lemma~\ref{lem:entropy-splitting}.  Retain typical parts
on which this cost and all finitely many score defects are bounded by
a common deterministic $q_\eps\to0$, using Markov thresholds fixed
before selecting parts.  The same lemma justifies the kernel
resampling with base $(T,C_r)$; at this stage no observation involving
a free replica is fixed.  Passing to finer finite meshes, and then to a diagonal,
gives the asserted common sequence and its dyadic shapes.  Shape bounds
follow from the fixed horizon.  A saturated observation at $u$,
localized with $r=d=0$ and normalized to duration one, has aspect
$e(u)/u$; minimality gives $e(u)\ge mu$.

Here are details of coherence.  Put
$D=|b(r)-b(t)|+|a(r)-a(t)|+t-r$.
Convert the observation at $t$ to the shape $(b(r),a(r))$ at $r$, using
its recorded frame rounded to the target precision and recording the
actual target cells.  Its additional conditional log-cardinality is at
most $C D+o(1)$.  To check this bound, rotate the old uncertainty into
the new frame: the tangent uncertainty entering the normal coordinate
has exponent at least
$b(t)+e(r)\ge a(r)-|b(t)-b(r)|$.  The other side-length changes cost
at most the absolute differences of the two depths.  Transport from
$C_t$ to $C_r$ adds $\eps^r$ times the velocity uncertainty, which has
the same bound in position units divided by $\eps^r$.  Counting the
finer cells in each resulting rectangle gives the stated estimate.
The weight changes by at most $2D$, and $S_t-S_r\le t-r+o(1)$.
Thus this converted test and the original test at $r$ have total
deficit at most $C_\gamma D+o(1)$ at duration $r$.

Call these two labels $D_1,D_2$, their common label $D_c$, and their
intersection label $D_i$.  They contain $C_r$, and their remaining
coordinates are deterministic functions of $T,C_r$ and their respective
orientations.  The resampling construction therefore gives
\[
 \Mut(D_1;D_2\mid T)\le S_r+o(1),\qquad
 \Ent(D_c\mid T)\ge S_r.
\]
The pair information identity and common-label determination imply
\begin{align*}
 \Mut(T;D_i)
 &\le\Mut(T;D_1)+\Mut(T;D_2)-\Ent(D_c)
                   +\Mut(D_1;D_2\mid T)+O(\xi)+o(1)\\
 &\le\Mut(T;D_1)+\Mut(T;D_2)-\Mut(T;D_c)+O(\xi)+o(1).
\end{align*}
Both $D_c,D_i$ have the same time $C_r$ and obey the original horizon.
Apply the universal score bound to each.  Their two time contributions
cancel the corresponding contributions in $D_1,D_2$, giving
\[
 2\gamma\EE(e(r)-v)\le C_\gamma D+O(\xi)+o(1).
\]
This inequality also holds on any separation event of fixed positive
limiting probability: binary splitting preserves the two original
saturations, and the extra conditional dependence costs only $o(1)$.
On the event $v\le d$ its left side is at least
$2\gamma(e(r)-d)$.  Choose $K_\gamma$ larger than the preceding
constant divided by $2\gamma$, and then let the angle mesh tend to
zero.  A positive limiting probability of separation at a depth in
\eqref{eq:classical-coherence} is impossible.

Apply the same comparison to two independent replicas of a single
label given $(T,C_u)$.  With $D=0$ they agree at every depth
$e(u)-\nu$, $\nu>0$, with probability tending to one.  For each base
value, a coarse cell and its bounded neighbors then contain all but
$o(1)$ of the conditional orientation mass on average.  Choose such a
cell deterministically from the base.  The entropy bound obtained by
separating this event from its complement is $o(1)$ at depth
$e(u)-\nu$; refining the orientation costs at most $\nu+o(1)$.
Let $\nu\downarrow0$.  This proves the asserted zero conditional
orientation entropy.
\end{proof}

\subsection{Isotropic tangents or a matched aspect}

\begin{proposition}[Tangent reduction]
\label{prop:isotropic-reduction}
One of the following alternatives holds for the least aspect $m$.
\begin{enumerate}
\item If $m=0$, there is a saturated sequence of duration and horizon
one, with isotropic endpoint $(b,a)=(0,0)$ and homogeneous conditional
time profile
\[
 \Ent(C_u\mid T)=su+o(1),\qquad 0\le u\le1,
 \quad 0\le s\le1.
\]
\item If $m>0$, the saturated path can be represented by a single
orientation $R$ with depth-$mu$ prefixes at time $u$, for every
$0<u\le1$, and
$\Ent(R_{mu}\mid T,C_u)=o(1)$.  Its shapes satisfy
\[
 e(u)=mu,\qquad b(u)\text{ is Lipschitz},\qquad
                 0\le-b'(u)\le1+m\quad\text{a.e.}
\]
There is $\kappa>0$ such that its limiting conditional orientation
entropies obey
\begin{equation}
 \Ent(R_{m(u+h)}\mid T,C_u)\ge\kappa h,
                  \qquad 0<u<u+h<1.
 \label{eq:orientation-growth}
\end{equation}
The alternatives and the positive lower bound persist after homogeneous
splittings that preserve the saturated path.
\end{enumerate}
\end{proposition}
\begin{proof}
The nesting inequalities make $b$ nonincreasing and $a(u)+u$
nondecreasing; extend them to their continuity points.  In particular
$e(u)+u$ is nondecreasing, so $e$ has bounded variation and its singular
distributional derivative is nonnegative.  Choose a common differentiability
point $x\in(0,1)$ for these functions.
Take dyadic $r<x<t$ with $\eta=t-r\downarrow0$ and with both distances
from $x$ comparable to $\eta$.  Differentiability gives uniform linear
approximations throughout $[r,t]$, with error $o(\eta)$.  Choose a finite
dyadic chain
\[
 r=u_0<u_1<\cdots<u_N=t
\]
whose largest interval length is $o(\eta)$, and put
\[
 D_j=|b(u_{j+1})-b(u_j)|+|a(u_{j+1})-a(u_j)|+u_{j+1}-u_j,
 \qquad D_{\max}=\max_jD_j.
\]
The uniform linear approximations give
$D_{\max}\le C\max_j(u_{j+1}-u_j)+o(\eta)=o(\eta)$.
Here the chain may become longer as $\eta\downarrow0$, but it is finite
and fixed at each inner small-parameter limit.

Let $e_{\min}=\min_j e(u_j)$ and choose a positive
$\delta_\eta=o(\eta)$.  Set
\[
 d=\bigl(e_{\min}-K_\gamma D_{\max}-\delta_\eta\bigr)_+.
\]
When $d>0$, it lies strictly below the coherence depth for every adjacent
pair in the chain.  Equation~\eqref{eq:classical-coherence} therefore
makes all their frames agree at depth $d$, up to bounded neighboring
cells, with probability tending to one.  Transitivity costs at most a
fixed number depending on this chain of neighboring cells, and the union
of its finitely many exceptional events still has probability tending
to zero.  The frame at $r$ is determined by $(T,C_r)$ to zero normalized
entropy at this depth.  The entropy ambiguity bound hence gives the same
determination for the terminal depth-$d$ frame.  If $d=0$, no orientation
information is required.

The uniform linear approximation for $e$ also gives
\[
 e(t)-e_{\min}\le\max(e'(x),0)\eta+o(\eta).
\]
Since $e_{\min}\ge0$, the definition of $d$ yields, in either case,
\[
 e(t)-d\le e(t)-e_{\min}+K_\gamma D_{\max}+\delta_\eta
       \le\max(e'(x),0)\eta+o(\eta).
\]
All depths, the chain, and its angular mesh are chosen before the inner
limit.  The depth loss is $K_\gamma D_{\max}$ at this one common depth;
the number of links enters only the zero-cost angular ambiguity.

Localize the terminal observation at $r$, using its common depth $b(t)$
and this $d$.  Its local horizon is $\eta+e(t)-d$.  Saturation and zero
conditional orientation entropy give saturated conditional systems.
With small parameter $\rho=\eps^\eta$, their duration is one and their
aspect is at most $\max(e'(x),0)+o(1)$.  All local coordinate bounds are
uniform, and their horizons stay bounded.  Minimality of $m$ therefore
implies $m\le\max(e'(x),0)$.  If $m>0$, this yields $e'(x)\ge m$
almost everywhere.  Since the singular derivative of $e$ is
nonnegative,
\[
 e(1)-e(u)\ge\int_u^1 e'(v)\,dv\ge m(1-u).
\]
Together with $e(u)\ge mu$ and $e(1)=m$, this proves $e(u)=mu$.
The original nesting inequalities now give
\[
 0\le b(r)-b(t)\le(1+m)(t-r),
\]
which proves the Lipschitz assertion and extends the path to every time.

We specify the selection of shrinking-scale limits, also for later use.
For the $j$th interval first fix its endpoints, its finite coherence
chain, and the first $j$ time-prefix fractions.  Homogenize their joint
and conditional probability bins, then choose the inner precision so
that all depth errors, normalized conditional entropies, uniform-bound
errors, and retained score deficits are less than $\eta_j/j^2$.
Use uniform typical sets before any further restriction.
The sum of the conditional deficits is then $o(\eta_j)$ in mean;
select typical localization cells with each required error
$o(\eta_j)$.  Only afterwards divide by $\eta_j$ and take the outer
diagonal.  This is an application of
Lemma~\ref{lem:profile-diagonal} with the error measured in the new
units.  No modulus of differentiability uniform over profiles is used.

 If $m=0$, use a common regular point with $e'\le0$ in the tangent
 construction above.  The set of such points has positive measure:
 otherwise $e'>0$ almost everywhere, and its nonnegative singular
 derivative would force $e(1)>e(u)\ge0$, contrary to $e(1)=0$.
 The nondecreasing $1$-Lipschitz homogeneous time profile $S_u$ is
 differentiable almost everywhere, so the point can also be chosen
 regular for $S$.  Put $s=S'(x)\in[0,1]$.  At this point the residual
 aspect just constructed is $o(\eta)$.

For the conditional time calculation, fix the finite list
$q=r+k\eta$ chosen above, and put $B=(T,C_r)$, $D_q=C_q$,
$d_q=S_q-S_r$, and $\ell=\log(1/\eps)$.  Let $\mathsf P$ be their
joint law with the localization label $J$, before its value is fixed.
The recorded conditional bins have one error $\omega_\eps=o(\eta)$
for this list outside probability $o(1)$, while the preceding parameter
choice gives $\Ent^{\mathsf P}(J\mid B)=o(\eta)$.  At sampled values,
\[
 -\ell^{-1}\log\mathsf P(D_q\mid B,J)
 +\ell^{-1}\log\mathsf P(D_q\mid B)
 =\ell^{-1}\log
       \frac{\mathsf P(J\mid B)}{\mathsf P(J\mid B,D_q)}.
\]
The two nonnegative $J$ surprises on the right have means at most
$\Ent^{\mathsf P}(J\mid B)$.  Choose $\delta_\eps=o(\eta)$ with
$\Ent^{\mathsf P}(J\mid B)/\delta_\eps\to0$ and trim them by Markov's
inequality.  Intersecting with the recorded bins leaves an event $E$ of
probability $1-o(1)$ on which, simultaneously for the listed $q$,
\[
 -\ell^{-1}\log\mathsf P(J\mid B)\le\delta_\eps,\qquad
 \left|-\ell^{-1}\log\mathsf P(D_q\mid B,J)-d_q\right|
       \le\alpha_\eps:=\omega_\eps+\delta_\eps=o(\eta).
\]
The list is fixed before choosing the inner precision; its errors and
exceptional masses are put under one deterministic vanishing envelope
before the lists grow along the preceding diagonal.

Put $\mathsf P^E=\mathsf P(\,\cdot\mid E)$.  Conditional relative
entropy gives
\[
 \Ent^{\mathsf P^E}(J\mid B)
 \le\mathbb E_{\mathsf P^E}[-\ell^{-1}\log\mathsf P(J\mid B)]
 \le\delta_\eps=o(\eta).
\]
Since $J$ records $R_d$, this also bounds the orientation cost under
$\mathsf P^E$.  The binary tag $\boldsymbol1_E$ costs
$O(1/\ell)=o(\eta)$.  Lemma~\ref{lem:classical-splitting}, the original
$o(\eta)$ deficits, and the common lower bound $-o(\eta)$ for the
complementary deficits give $o(\eta)$ deficits on $\mathsf P^E$.
Localization of this law has conditional laws
$\mathsf Q_j=\mathsf P(\,\cdot\mid E,J=j)$, followed by the known
$J$-dependent coordinate change.  The shrinking-scale selection above
therefore gives a set of $\mathsf P^E_J$-mass $1-o(1)$ on which the
required local deficits are $o(\eta)$ under the same finite envelope.

Choose one $\lambda_\eps\downarrow0$ with
$\log(1/\lambda_\eps)=o(\eta\ell)$, using $\eta\ell\to\infty$.
Apply Lemma~\ref{lem:typical-conditioning} to the event $E$, with
$C=J$, $X=(B,J)$, $Y=D_q$, and all three exceptional thresholds
equal to $\lambda_\eps$.
Outside $\mathsf P^E_J$-mass at most $\lambda_\eps$, it gives
\[
 \mathsf Q_j\le[\mathsf P(E)\lambda_\eps]^{-1}\mathsf P_j,
 \qquad \mathsf P_j=\mathsf P(\,\cdot\mid J=j).
\]
On $J=j$, conditioning on $X$ is conditioning on $B$, and the bin
on $E$ is for $\mathsf P(D_q\mid B,J)=\mathsf P_j(D_q\mid B)$.
The lemma's entropy and surprise conclusions therefore give
\[
 \begin{gathered}
 \Omega_\eps:=
 \alpha_\eps+\ell^{-1}\log\frac1{\mathsf P(E)\lambda_\eps^2}
       =o(\eta),\\
 \left|\Ent^{\mathsf Q_j}(D_q\mid B)-d_q\right|\le\Omega_\eps,\qquad
 \mathsf Q_j\!\left\{
 \left|-\ell^{-1}\log\mathsf Q_j(D_q\mid B)-d_q\right|
       >\Omega_\eps\right\}\le2\lambda_\eps .
 \end{gathered}
\]
These are uniform over the retained $j$ and the fixed finite list;
the finite union bound is absorbed by the same diagonal.  No original
cell probability enters the error envelope.

Intersect these cells with the typical localization cells above.
Since $C_r$ is fixed by $J$ and the local time grids correspond to
$C_q$ up to bounded ambiguity, the estimates are conditional on the
local particle.  Differentiability gives $(S_q-S_r)/\eta=ks+o(1)$.
The residual aspect and excess horizon are $o(1)$ in the new units;
coarsen these vanishing depth errors and apply the universal bound to
recover a saturated isotropic sequence.  Increasing the finite prefix
lists gives the homogeneous linear time profile and proves the first
alternative.

Suppose henceforth $m>0$.  Fix a dyadic $u>0$.  On a sufficiently fine
finite chain from $u$ to $1$, coherence and the Lipschitz width bounds
identify all frames at every depth strictly below $mu$.  The terminal
orientation $R$ thus determines the frame at $u$ at that precision.
Let the gap from $mu$ tend to zero.  Both the old and new rectangles
then determine each other with vanishing log-cardinality error;
recording the same data in the endpoint-prefix frame preserves its
score.  Taking a diagonal over dyadic $u$, and then using Lipschitz
interpolation of the widths and time grids, supplies the common-prefix
representation at every $u$.

Finally, if no $\kappa>0$ satisfies \eqref{eq:orientation-growth},
there are fixed intervals $(u_j,u_j+h_j)$, chosen after their inner
limits, with
\[
 h_j^{-1}\Ent(R_{m(u_j+h_j)}\mid T,C_{u_j})\longrightarrow0.
\]
Apply Lemma~\ref{lem:classical-localization} to the saturated observation
at $u_j+h_j$, with $r=u_j$, $b=b(u_j+h_j)$, and
$d=e(u_j+h_j)$.  Its conditional final shape is exactly isotropic and
its local horizon is $h_j$.  The error $2\zeta$ in
\eqref{eq:classical-local-first} is $o(h_j)$.  The other composition
errors are made $o(h_j)$ by first fixing the interval and then choosing
the inner precision.  The local deficits consequently have mean
$o(h_j)$ and a uniform lower bound of the same order.  Select a local
system with deficit $o(h_j)$ and replace $\eps$ by $\eps^{h_j}$.
This produces a saturated isotropic sequence of horizon one, contrary
to $m>0$.  This contradiction proves the uniform positive constant.
It uses neither a positive lower bound for $h_j$ nor one for
$u_j$: both are fixed positive numbers at their own inner stage.

The same arguments apply after any homogeneous splitting preserving
the path's scores.  Such a path is still a least-aspect extremizer, so
the final contradiction again supplies a positive constant.  One can
also fix a projective orientation chart by a finite split.  In that
chart use rows $(1,0)$ and $(-R,1)$, with bounded slope $R$ and nested
slope prefixes.  The row frame and orthogonal frame have bounded
condition number, angular and slope precisions are comparable, and
the resulting rectangle observations differ by bounded covering
multiplicities.  All conclusions therefore hold in this row convention.
\end{proof}

\begin{remark}[Homogeneous profiles for subsequent local arguments]
\label{rem:classical-profile-inheritance}
In either alternative of Proposition~\ref{prop:isotropic-reduction},
one may first impose a fixed finite list of pairs $(X_j,Y_j)$ satisfying
$|\mathcal Y_j|\le\eps^{-A_j}$ for fixed $0\le A_j<\infty$ and all
sufficiently small $\eps$.  The children may be joint grid observations
or geometric prefixes.  No cardinality bound is imposed on the
conditioning variables $X_j$, which may include $T$.
  Apply
Lemma~\ref{lem:homogenization} and
Lemma~\ref{lem:classical-splitting}, select a saturated homogeneous
component, and only then define its limiting profiles and choose a
differentiability or trace point.  The profiles need not equal those
before this initial splitting.  The positive constant in
\eqref{eq:orientation-growth} is obtained for the selected component
by the same least-aspect contradiction.

To spell out the later inheritance, suppose homogeneous probability
bins for a parent $P$ and a refinement $Q$ have exponents $d_P,d_Q$
with error $\eta$; both may be conditional on an additional base $W$,
including the exact particle.
Here $d_P$ is the exponent for the child $P$, and $d_Q$ is the exponent
for the joint child $(P,Q)$, identified with $Q$ when $Q$ already records
$P$; both children are conditional on $W$.  If the exact particle is
present, it belongs to $W$.
  On their common uniform typical set,
ratios of joint and parent masses give
\[
 \Pr(Q=q\mid P=p,W=w)
       \le\eps^{d_Q-d_P-2\eta},\qquad
 \#\{q:\ (p,q,w)\text{ retained}\}
       \le\eps^{-(d_Q-d_P)-2\eta}.
\]
These are respectively an upper conditional mass and an upper retained
support count; they are stronger than an averaged entropy inequality.
If a subsequent normalized law $\nu$ obeys
$\nu\le\eps^{-\theta}\mu$, discard parent values whose marginal
likelihood $\nu_{P,W}/\mu_{P,W}$ is below $\eps^\lambda$.  Their
$\nu$-probability is at most $\eps^\lambda$, and on the remaining
parents the conditional mass bound loses at most $\theta+\lambda$
in its exponent.  Support counts are inherited under restriction.
The conditional relative-entropy identity in
Lemma~\ref{lem:homogenization} supplies the corresponding lower entropy
bound with loss $\theta$.  There is no factor depending on the number
of parent values.  Repeat this operation for the fixed finite list,
then take a diagonal.  For a block of length $h$, choose
$\eta,\theta,\lambda=o(h)$ at its inner stage before using
$\rho=\eps^h$ as the small parameter.  Cap regularization may be imposed
at the same stage to retain the root bound
\eqref{eq:classical-cap-bounds} under these operations.
\end{remark}
